\section{Conclusions}
\label{sec:conclusions}

Across the laboratory and field studies reviewed here, the response to a sound depends as much on the receiver as on the sound. Neither discipline explains this alone. Field studies show what natural signals do under real costs; laboratory studies show how the brain turns sound into physiology and action, mostly with sounds no rodent hears in the wild.

Rodent alarm systems carry several kinds of information. Marmot call rate tracks risk \citep{BlumsteinArmitage1997}, Gunnison's prairie dog calls differ with the type of threat \citep{KiriazisSlobodchikoff2006}, gray squirrels respond to unfamiliar heterospecific alarms that resemble familiar ones \citep{GetschowEtAl2013}, and after a hawk call, ordinary bird chatter lets squirrels relax sooner than ambient noise does \citep{LillyEtAl2019}. Yet natural calls are less special than they first appear. Once controls share a call's frequency band and timing, the natural call's advantage shrinks: pure tones still draw rats toward the speaker, and frequency-matched tones reproduce most of the amygdala's discrimination between call types \citep{WohrSchwarting2007,ParsanaEtAl2012a}. The informative question is which features a receiver uses, not whether a sound is natural.

Recognizing a sound and acting on it are separable. Naïve rats slow their movement and heart rate to 22-kHz alarm calls but freeze selectively only after aversive experience \citep{OlszynskiEtAl2020,KimEtAl2010}, and juvenile ground squirrels discriminate alarm types by heart rate before they do so behaviorally \citep{Mateo1996}. Mice hearing the same sweep freeze or flee depending on their surroundings \citep{MongeauEtAl2003}, and chipmunks hearing alarms trade vigilance against foraging and distance to cover \citep{BaackSwitzer2000}. A behavioral response is therefore not a direct measure of how threatening a sound is.

Experience reshapes responding in ways a single decline curve cannot reveal. Juvenile ground squirrels discount callers whose alarms prove unreliable, whereas marmots respond more to them \citep{HareAtkins2001,BlumsteinEtAl2004}. Responses that wane to one caller return when the caller changes \citep{Hare1998}, a result that fits stimulus-specific habituation as well as learning about the caller. Irregular timing slows habituation, but so far only to artificial sounds \citep{Davis1970,HerryEtAl2007}. Telling these processes apart requires measuring recovery, source specificity, timing, and costs in the same animals.

The neural evidence is the thinnest link. Most causal circuit maps were built with tones, noise, and looming sounds; for a natural alarm call, the clearest causal result remains the requirement for the auditory thalamus in experience-gated freezing to 22-kHz calls \citep{KimEtAl2010}. In sciurids, natural alarm calls have been linked to region-specific neural activation \citep{FreemanEtAl2020} but not, to our knowledge, to a causally established circuit. The receiver-centered framework offered here---rapid acoustic defense, learned threat evaluation, and flexible action selection---is a way to work across that gap: for any finding, it asks which demand the finding bears on, which prediction would test it (\tabref{tab:predictions}), and how far it can be carried between laboratory and field.

A sound does not dictate an animal's response. Its effects depend on what the receiver can hear, what it has learned, what it stands to gain or lose, and what has happened before. Field studies reveal how these factors shape behavior under natural conditions. Laboratory neuroscience identifies mechanisms through which acoustic information can influence learning, physiology, and action, although their roles in processing natural calls remain incompletely understood. Bringing these perspectives together allows us to ask more precise questions about how biological significance emerges---and why the same sound can lead to different outcomes.
